\documentclass[11pt]{article}
\usepackage[margin=1in]{geometry}
\usepackage{amsmath,amssymb,amsthm}
\usepackage{graphicx}
\usepackage[hidelinks]{hyperref}

\newtheorem{theorem}{Theorem}
\newtheorem{lemma}{Lemma}
\newtheorem{remark}{Remark}

\newcommand{\T}{\mathbb T}
\newcommand{\R}{\mathbb R}
\newcommand{\Z}{\mathbb Z}
\newcommand{\cP}{\mathcal P}
\newcommand{\cS}{\mathcal S}
\newcommand{\norm}[1]{\left\lVert #1 \right\rVert}
\newcommand{\Decomp}{\mathsf D}

\title{Torus and Real-Line Fourier Decomposition Properties Are Equivalent}
\author{FA Banach run \texttt{fa\_banach\_001}, agent \texttt{agent\_lane\_05}}
\date{2026-06-16}

\begin{document}
\maketitle

\section*{Status}

This packet records a candidate full solution, likely valid, to Problem 5.2 of
Deng, Lorist, and Veraar \cite{DLV}: whether the Fourier decomposition
properties on \(\T\) are equivalent to their counterparts on \(\R\).

\section*{Source Passage}

\begin{center}
\includegraphics[width=0.98\linewidth]{figures/open_problem_crop.png}
\end{center}

\section*{Transcription of the Open Problem}

The source paper writes:
\begin{quote}
``Our decomposition properties are Fourier decomposition properties on
\(\T\). One may similarly define Fourier decomposition properties on
\(\R\), in which case it is natural to wonder if these properties would be
equivalent.''
\end{quote}
Problem 5.2 asks:
\begin{quote}
``Are the decomposition properties equivalent to their counterparts on
\(\R\)?''
\end{quote}

\section*{Definitions Used}

Let \(G=\T\) or \(G=\R\), and let \(\widehat G=\Z\) or \(\widehat G=\R\),
respectively.  For an interval \(I\subseteq \widehat G\), write
\[
  D_I^G f = \mathcal F_G^{-1}({\bf 1}_I \widehat f).
\]
For a finite interval family \(\mathcal I\), put
\[
  A_{\mathcal I}^G f = (D_I^G f)_{I\in\mathcal I}.
\]
The target norm is the mixed norm
\[
  \norm{(g_I)_{I\in\mathcal I}}_{\ell^q(L^p(G;X))}
  =
  \begin{cases}
  \left(\sum_{I\in\mathcal I}\norm{g_I}_{L^p(G;X)}^q\right)^{1/q},
    & q<\infty,\\
  \sup_{I\in\mathcal I}\norm{g_I}_{L^p(G;X)},& q=\infty.
  \end{cases}
\]
Lower \(\ell^q(L^p)\)-decompositions mean that
\[
  \norm{A_{\mathcal I}^G f}_{\ell^q(L^p(G;X))}
  \le L\norm{f}_{L^p(G;X)}
\]
uniformly over interval partitions and test functions.  Upper
\(\ell^q(L^p)\)-decompositions mean the reverse estimate
\[
  \norm{f}_{L^p(G;X)}
  \le U\norm{A_{\mathcal I}^G f}_{\ell^q(L^p(G;X))}
\]
uniformly over interval partitions.  As usual, lower estimates may be checked
on finite subfamilies.  Upper estimates may be checked on finite subfamilies
which contain the Fourier support of the test function, since the remaining
partition pieces have zero projection.  One then passes to general
trigonometric polynomials on \(\T\), or Schwartz functions with compact
Fourier support on \(\R\), by density and monotone convergence.

\section*{Theorem}

\begin{theorem}
Let \(X\) be a Banach space, let \(1<p<\infty\), and let
\(1\le q\le\infty\).  Then \(X\) has lower \(\ell^q(L^p)\)-decompositions on
\(\T\) if and only if it has lower \(\ell^q(L^p)\)-decompositions on \(\R\).
The best constants are equal.

The same equivalence, again with the same best constants, holds for upper
\(\ell^q(L^p)\)-decompositions.
\end{theorem}

\section*{Proof Intuition}

The decomposition inequalities only involve finitely many interval projections
when tested on a polynomial or on a Schwartz function with compact Fourier
support.  For such a finite family, the map \(f\mapsto(D_I f)_I\) is just a
finite tuple of Fourier multipliers.  De Leeuw's periodization theorem
compares exactly these multipliers on \(\R\) with their sampled versions on
\(\T\), and de Leeuw's restriction theorem compares a real-line multiplier
with its values on integer frequencies.

The apparent obstacle is that the target norm is not a pointwise
\(L^p(G;\ell^q(X))\) norm, but an \(\ell^q\)-sum of \(L^p\) norms.  This causes
no loss for finite interval families: the standard periodization and
restriction proofs give convergence of every component in \(L^p\), and a
finite \(\ell^q\) norm commutes with taking the limit.  Thus both sides of the
decomposition inequalities transfer unchanged.

\begin{lemma}[Finite mixed-norm de Leeuw principle]\label{lem:deleeuw}
Let \(1<p<\infty\), \(1\le q\le\infty\), and let \(X\) be a Banach space.

\begin{enumerate}
\item Let \(\mathcal J=\{J_1,\ldots,J_N\}\) be a finite family of disjoint
intervals in \(\R\), and let \(f\in \cS(\R;X)\) have compact Fourier support
whose boundary points do not meet the endpoints of the \(J_j\)'s.  There are
numbers \(\varepsilon_n\downarrow0\), trigonometric polynomials \(f_n\), and
finite interval families \(\mathcal J_n\) in \(\Z\) such that
\[
  \norm{f_n}_{L^p(\T;X)}\to\norm{f}_{L^p(\R;X)}
\]
and, for each \(j\),
\[
  \norm{D_{J_{j,n}}^\T f_n}_{L^p(\T;X)}
  \to
  \norm{D_{J_j}^\R f}_{L^p(\R;X)}.
\]

\item Let \(\mathcal I=\{I_1,\ldots,I_N\}\) be a finite family of disjoint
intervals in \(\Z\), and let \(f\in\cP(\T;X)\).  There are Schwartz functions
\(F_n\in \cS(\R;X)\) and finite interval families \(\mathcal I_n\) in \(\R\)
such that
\[
  \norm{F_n}_{L^p(\R;X)}\to\norm{f}_{L^p(\T;X)}
\]
and, for each \(j\),
\[
  \norm{D_{I_{j,n}}^\R F_n}_{L^p(\R;X)}
  \to
  \norm{D_{I_j}^\T f}_{L^p(\T;X)}.
\]
\end{enumerate}
Consequently the same convergence holds for the finite
\(\ell^q(L^p)\)-norms of the tuples of interval projections.
\end{lemma}

\begin{proof}
This is the usual proof of de Leeuw's compactification and restriction
theorems for \(L^p\)-multipliers, applied to the finite tuple of symbols
\(({\bf 1}_{J_j})_{j=1}^N\) or \(({\bf 1}_{I_j})_{j=1}^N\).  We recall the
points relevant to the mixed norm.

For compactification, sample the real-line multiplier at the lattice
\(\varepsilon\Z\).  The preimages
\[
  J_{j,\varepsilon}:=\{k\in\Z:\varepsilon k\in J_j\}
\]
are intervals in \(\Z\), and they are disjoint when the \(J_j\)'s are
disjoint.  The standard periodization argument, with the usual normalization
of dilates, gives
\[
  \norm{f_\varepsilon}_{L^p(\T;X)}
  \to \norm{f}_{L^p(\R;X)}
\]
and, for each fixed \(j\),
\[
  \norm{D_{J_{j,\varepsilon}}^\T f_\varepsilon}_{L^p(\T;X)}
  \to
  \norm{D_{J_j}^\R f}_{L^p(\R;X)}.
\]
The endpoint-avoidance assumption on \(\widehat f\) removes the only
discontinuity issue; a standard approximation removes it for general test
functions.

For restriction, choose real intervals \(\widetilde I_j\subset\R\) with
\(\widetilde I_j\cap\Z=I_j\), for instance by putting endpoints halfway
between adjacent integers.  The classical restriction proof localizes a
rescaled copy of the trigonometric polynomial on \(\R\) so that its Fourier
mass concentrates near the integer frequencies.  Applying this to each
\(\widetilde I_j\) yields the displayed componentwise convergence.

These are exactly de Leeuw's arguments \cite{deLeeuw}; see also the
operator-valued presentation in \cite[Chapter 5]{HNVW1}.  Since the family has
only \(N<\infty\) components, componentwise convergence of the \(L^p\)-norms
implies convergence of the finite \(\ell^q\)-norm, including \(q=\infty\).
\end{proof}

\begin{proof}[Proof of the theorem]
We first prove the lower estimate.  Suppose \(X\) has lower
\(\ell^q(L^p)\)-decompositions on \(\T\) with constant \(L\).  Let
\(\mathcal J=\{J_1,\ldots,J_N\}\) be a finite disjoint interval family in
\(\R\), and let \(f\in\cS(\R;X)\) have compact Fourier support.  By a harmless
approximation we may assume that \(\widehat f\) does not charge the endpoints
of the \(J_j\)'s.  Apply Lemma \ref{lem:deleeuw}(1).  For every \(n\),
\[
  \norm{A_{\mathcal J_n}^\T f_n}_{\ell^q(L^p(\T;X))}
  \le L\norm{f_n}_{L^p(\T;X)}.
\]
Passing to the limit gives
\[
  \norm{A_{\mathcal J}^\R f}_{\ell^q(L^p(\R;X))}
  \le L\norm{f}_{L^p(\R;X)}.
\]
Density and monotone convergence over finite subfamilies give the real-line
lower decomposition property.

Conversely, suppose the lower estimate holds on \(\R\) with constant \(L\).
Let \(\mathcal I=\{I_1,\ldots,I_N\}\) be a finite disjoint interval family in
\(\Z\), and let \(f\in\cP(\T;X)\).  Lemma \ref{lem:deleeuw}(2) gives
real-line approximants \(F_n\) and interval families \(\mathcal I_n\).  The
real-line estimate gives
\[
  \norm{A_{\mathcal I_n}^\R F_n}_{\ell^q(L^p(\R;X))}
  \le L\norm{F_n}_{L^p(\R;X)}.
\]
Passing to the limit yields the corresponding estimate on \(\T\).  Thus the
best lower constants on \(\T\) and \(\R\) are equal.

For upper decompositions, first take a finite interval family containing the
Fourier support of the test function; equivalently, take the finite part of an
interval partition on which the test function has nonzero projections.  The
same approximation lemma preserves this support condition: the preimage
intervals contain the Fourier support of the approximating trigonometric
polynomial, and in the restriction direction the thickened real intervals
contain the Fourier support of the real-line approximants.  The proof is then
identical, with the inequality reversed.  If the upper estimate holds on
\(\T\), then for the approximants in Lemma \ref{lem:deleeuw}(1),
\[
  \norm{f_n}_{L^p(\T;X)}
  \le U\norm{A_{\mathcal J_n}^\T f_n}_{\ell^q(L^p(\T;X))},
\]
and the limit gives the real-line upper estimate.  If the upper estimate
holds on \(\R\), Lemma \ref{lem:deleeuw}(2) transfers it back to \(\T\).  The
best upper constants are therefore equal as well.
\end{proof}

\section*{Verification Notes}

No numerical verification is involved.  The proof reduces the open problem to
the standard de Leeuw compactification/restriction limit and checks that the
finite mixed norm used in Deng--Lorist--Veraar's decomposition definition is
preserved under that limit.  The argument is insensitive to the precise
choice of open, closed, or half-open interval endpoints, because the estimates
are first proved for test functions whose Fourier supports avoid endpoints
and then extended by approximation.

\section*{Novelty and Literature Check}

The run indexes were searched for arXiv:2311.01436, ``Strongly Kreiss'',
``UMD'', ``Fourier decomposition'', ``decomposition properties on R'', and
related torus/real-line transference phrases.  No existing packet or attempt
addresses this problem.

External web searches for exact phrases around the paper title, Problem 5.2,
``ell\^{}q(L\^{}p)-decompositions'', and torus/real-line decomposition
transference found the source paper and adjacent multiplier literature, but
no later paper explicitly resolving this open problem.
Novelty confidence is moderate: the proof is a direct application of standard
de Leeuw machinery once the mixed-norm finite-family formulation is isolated.

\section*{Limitations}

The theorem is stated in the nontrivial \(L^p\) range \(1<p<\infty\), which is
the range used in the source paper's UMD/Fourier-multiplier theory.  Endpoint
conventions for \(p=1\) or \(p=\infty\) are not addressed here.  The packet
does not solve the scalar endpoint Problem 5.3 or the sharp extrapolation
Problem 5.4.

\section*{Human Review Recommendation}

Reviewers should check Lemma \ref{lem:deleeuw}: in particular, that the
standard de Leeuw compactification and restriction proofs may be applied
componentwise to finite interval families and that the convergence of
component \(L^p\)-norms is enough to pass to the finite
\(\ell^q(L^p)\)-mixed norm without a cardinality-dependent loss.

\begin{thebibliography}{9}

\bibitem{DLV}
C. Deng, E. Lorist, and M. C. Veraar,
\emph{Strongly Kreiss Bounded Operators in UMD Banach Spaces},
arXiv:2311.01436.

\bibitem{deLeeuw}
K. de Leeuw,
\emph{On \(L_p\) multipliers},
Ann. of Math. (2) \textbf{81} (1965), 364--379.

\bibitem{HNVW1}
T. Hyt\"onen, J. van Neerven, M. Veraar, and L. Weis,
\emph{Analysis in Banach Spaces. Volume I: Martingales and Littlewood--Paley
Theory},
Ergebnisse der Mathematik und ihrer Grenzgebiete 63, Springer, 2016.

\end{thebibliography}

\end{document}
